% source: https://github.com/pfradin/luadraw/discussions/151#discussioncomment-15003356 (pfradin, block 1)
\documentclass[varwidth,border = 3pt]{standalone}
\usepackage[svgnames]{xcolor}
\usepackage{luadraw}
\usepackage{fouriernc}

\begin{luacode*}
function solve(f,a,b,n,df)
    n = n or 25
    if a > b then a, b = b, a end
    local x, delta, S = a, (b-a)/n, {}
    local fin, h, r = x+delta, 1e-6, nil
    if df == nil then
        df = function(x)
            return (f(x+h)-f(x))/h
        end
    end
    local iter = function(x)
        return x-f(x)/df(x)
    end
    while x < b do
        r = x+delta/2
        for i = 1,5 do
            if r~= nil then r = evalf(iter,r) end
            if (r==nil) then break end
        end
        if (r ~= nil) and (r >= x) and (r < fin) and (math.abs(f(r)) < 1e-6) then
            table.insert(S,r) -- on considère que r est une solution
        end
        x = fin; fin = x+delta
    end
    if #S > 0 then return S end
end

function tangent_from(p,t1,t2,from,dp)
-- p: t -> p(t) parameterization of the curve
-- t1, t2 : interval boundaries
-- from : point on the plane from which the tangents will originate
-- dp (optionnal) derivative of the function p
    local h = 1e-6
    if dp == nil then
        dp = function(t) -- approximate derivative function
            return (p(t+h)-p(t))/h
        end
    end
    local f = function(t)
        return cpx.det(p(t)-from,dp(t))
    end
    local n = math.max(25, math.floor((t2-t1)*2.5))
    local S = solve(f,t1,t2-h,n) --<< here we use the solve function
    if S ~= nil then return map(p,S) else return {} end -- empty list if no solution is found
end
\end{luacode*}

\begin{document}

\begin{luacode*}
function draw(index)
local g = graph:new{bbox=false, size={10,10}}
local i = cpx.I
local cos, sin, pi = math.cos, math.sin, math.pi

local p1 = function(t)
    return t+i*t^2/4 -- Explorer example
end
local p2 = function(t)
    return  1+i+rotate( 2*cos(t)+i*sin(t),-15)  -- an ellipse
end
local p3 = function(t)
    return  t + i*(t+1)*(t-2)*t
end
local L = {p1,p2,p3}
local p = L[index]
local P = 1.5-1.5*i -- the "from" point
local S = tangent_from(p,-5,5,P)

-- drawing
local axis_style = {
  arrows='-Stealth',
  xyticks={0.15, 0.15},
  legend={'$x$','$y$'} 
}
g:Daxes({0, 1, 1}, axis_style)
g:Dparametric(p,{t={-5,5}, draw_options="line width=0.8pt, red"})
for _,a in ipairs(S) do
    g:Dline(P,a,"blue")
end
g:Ddots(S,"red"); g:Ddots(P); g:Dlabel("$P$",P,{pos="S"})
g:Sendtotex()
end
\end{luacode*}

\centering

\section*{Explorer example}

\directlua{draw(1)}%

\section*{Another example}

\directlua{draw(2)}%

\section*{Not a quadratic curve}

\directlua{draw(3)}%

\end{document}
